% TOP_PROBLEM: {"id": 30006865, "title": "Arnold Chord Conjecture in Higher Dimensions", "category_id": 6, "category": {"id": 6, "name": "geometry", "display_name": "Geometry", "description": "Euclidean and non-Euclidean geometry, geometric structures.", "slug": "geometry", "order_index": 6, "created_at": "2026-07-31T15:26:25.671Z"}, "status": "open"}
% FIRST_POSTED: 2026-09-25
% ENTRY_KIND: statement_only
% !TeX program = lualatex
% Definition and sources only; this file is not an LLM solution attempt.
\documentclass[11pt]{article}
\usepackage[margin=25mm]{geometry}
\usepackage{fontspec}
\setmainfont{Latin Modern Roman}
\usepackage{amsmath,amssymb,mathtools}
\usepackage{unicode-math}
\setmathfont{Latin Modern Math}
\usepackage{xurl,hyperref}
\hypersetup{colorlinks=true,urlcolor=blue}
\setcounter{secnumdepth}{0}
\setlength{\parindent}{0pt}
\setlength{\parskip}{5pt}
\setlength{\emergencystretch}{3em}
\begin{document}
\section*{\texorpdfstring{Arnold Chord Conjecture in Higher Dimensions}{Arnold Chord Conjecture in Higher Dimensions}}\label{30006865}
\subsection{Definitions and mathematical statement}
A co-oriented contact structure on a closed $(2n+1)$-manifold $Y$ is $\xi=\ker\alpha$ for a one-form with $\alpha\wedge({\,\mathrm{d}}\alpha)^n$ nowhere zero. Its Reeb vector field is uniquely determined by
\[
 \alpha(R_\alpha)=1,\qquad \iota_{R_\alpha}{\,\mathrm{d}}\alpha=0.
\]
A Legendrian submanifold $\Lambda\subset Y$ is a closed $n$-dimensional submanifold with $T\Lambda\subset\xi$. The higher-dimensional chord assertion says that for every such datum with $n\ge2$, there are $T>0$ and $x\in\Lambda$ such that $\varphi_{R_\alpha}^T(x)\in\Lambda$. Thus $t\mapsto\varphi_{R_\alpha}^t(x)$, $0\le t\le T$, is a nonconstant Reeb chord. Its two endpoints need not coincide, and a closed Reeb orbit disjoint from $\Lambda$ would not answer the question.
\par\smallskip\textit{Formulation and concept references:} \href{https://arxiv.org/abs/1004.4319}{[S1]}.

\subsection{Short English statement}
Must every closed Legendrian submanifold in a closed higher-dimensional contact manifold have a nonconstant Reeb trajectory beginning and ending on it?
\subsection{Sources}
\begin{itemize}
\item[S1]Proof of the Arnold chord conjecture in three dimensions I (Hutchings-Taubes, arXiv abstract). Accessed 2026-09-01.
\url{https://arxiv.org/abs/1004.4319}.
\textit{Formulation source.}
\end{itemize}
\end{document}
